\documentclass[11pt]{article}
\usepackage[utf8]{inputenc}
\usepackage{graphicx}
\usepackage[margin=1in, includeheadfoot, footskip=133pt]{geometry}
\usepackage{xcolor}
\usepackage{hyperref}
\usepackage{fancyhdr}
\usepackage{enumitem}
\usepackage{titlesec}
\usepackage{amsmath}
\usepackage{amssymb}
\usepackage{tikz}
\usepackage{unicode-math}
\usetikzlibrary{arrows,arrows.meta,calc,positioning}

\newsavebox\pandoc@box
\newcommand*\pandocbounded[1]{% scales image to fit in text height/width
  \sbox\pandoc@box{#1}%
  \Gscale@div\@tempa{\textheight}{\dimexpr\ht\pandoc@box+\dp\pandoc@box\relax}%
  \Gscale@div\@tempb{\linewidth}{\wd\pandoc@box}%
  \ifdim\@tempb\p@<\@tempa\p@\let\@tempa\@tempb\fi% select the smaller of both
  \ifdim\@tempa\p@<\p@\scalebox{\@tempa}{\usebox\pandoc@box}%
  \else\usebox{\pandoc@box}%
  \fi%
}

% Define Caltech orange color
\definecolor{caltechorange}{RGB}{255, 108, 12}

% Define tightlist command
\providecommand{\tightlist}{%
  \setlength{\itemsep}{0pt}\setlength{\parskip}{0pt}}

% Section formatting
\titleformat{\section}
  {\normalfont\Huge\bfseries}{}{0em}{}
\titlespacing*{\section}
  {0pt}{3.5ex plus 1ex minus .2ex}{2.3ex plus .2ex}

% Subsection formatting - no numbers
\titleformat{\subsection}
  {\normalfont\Large\bfseries}{}{0em}{}
\titlespacing*{\subsection}
  {0pt}{3.25ex plus 1ex minus .2ex}{1.5ex plus .2ex}

% Subsubsection formatting - no numbers  
\titleformat{\subsubsection}
  {\normalfont\large\bfseries}{}{0em}{}
\titlespacing*{\subsubsection}
  {0pt}{3.25ex plus 1ex minus .2ex}{1.5ex plus .2ex}


% Page style
\pagestyle{fancy}
\fancyhf{}

% Footer setup
\cfoot{\thepage}

% Header setup
\fancyhead[L]{%
    \includegraphics[height=1cm]{templates/caltech-ctme-logo.png}
}
\fancyhead[R]{%
    {\color{gray}\small\href{https://ctme.caltech.edu}{ctme.caltech.edu}}
}

% Add padding after header
\setlength{\headheight}{50pt}
\setlength{\headsep}{2em}

% Force fancy style on first page
\AtBeginDocument{\thispagestyle{fancy}}



% List formatting
\setlist{nosep}  % Remove vertical spacing between list items
\setlist[itemize]{leftmargin=*}  % Align bullet points with left margin

% Hyperref setup
\hypersetup{
    colorlinks=true,
    linkcolor=caltechorange,
    filecolor=caltechorange,
    urlcolor=caltechorange,
    citecolor=caltechorange
}

% Document
\begin{document}

\section{Feature Scaling and Variable
Importance}\label{feature-scaling-and-variable-importance}

\subsection{1. Why Scale Features?}\label{why-scale-features}

\subsubsection{1.1 Motivation: A Logistic Regression
Example}\label{motivation-a-logistic-regression-example}

Consider predicting heart disease using two variables: - Age (years):
typically 20-80 - Cholesterol (mg/dL): typically 150-300

Without scaling:

\begin{verbatim}
# Simulated data
age = np.random.uniform(20, 80, 1000)        # age in years
chol = np.random.uniform(150, 300, 1000)     # cholesterol in mg/dL
# True relationship: equal importance
y = 1/(1 + np.exp(-(0.5*(age-50)/30 + 0.5*(chol-225)/75)))
y = (np.random.random(1000) < y).astype(int)

# Fit logistic regression
from sklearn.linear_model import LogisticRegression
model_raw = LogisticRegression()
model_raw.fit(np.column_stack([age, chol]), y)
print("Raw coefficients:", model_raw.coef_[0])
# Output: [-0.02, 0.005] # Cholesterol appears less important!
\end{verbatim}

With scaling:

\begin{verbatim}
# Scale features to zero mean, unit variance
age_scaled = (age - 50)/30     # (x - mean)/std
chol_scaled = (chol - 225)/75  # (x - mean)/std

model_scaled = LogisticRegression()
model_scaled.fit(np.column_stack([age_scaled, chol_scaled]), y)
print("Scaled coefficients:", model_scaled.coef_[0])
# Output: [0.5, 0.5] # Now we see equal importance!
\end{verbatim}

Key insights: 1. Without scaling: cholesterol's coefficient looks tiny
2. With scaling: true relationship becomes clear 3. Both variables
equally affect disease risk 4. Scaling reveals actual variable
importance

\begin{figure}
\centering
\pandocbounded{\includegraphics[keepaspectratio]{templates/011/scaling_effect.png}}
\caption{Effect of scaling on logistic regression}
\end{figure}

The plots above show: - Left: Raw data with very different scales -
Right: Scaled data reveals true relationship - Colors indicate disease
presence/absence - Notice how decision boundary is clearer in scaled
space

\begin{figure}
\centering
\pandocbounded{\includegraphics[keepaspectratio]{templates/011/coefficient_comparison.png}}
\caption{Coefficient comparison}
\end{figure}

The bar plot shows: - Raw coefficients are hard to compare - Scaled
coefficients reveal equal importance - Without scaling, we might wrongly
ignore cholesterol

Notes: - The true relationship uses scaled variables internally (note
the /30 and /75) - Coefficients in logistic regression show log-odds
impact - Age range: 60 years (20-80) - Cholesterol range: 150 mg/dL
(150-300) - Without scaling, a 1-unit change in cholesterol (1 mg/dL)
appears less important than a 1-unit change in age (1 year) - After
scaling, both variables are in standard deviation units, making
coefficients directly comparable

\subsubsection{1.2 Common Scaling Methods}\label{common-scaling-methods}

\begin{enumerate}
\def\labelenumi{\arabic{enumi}.}
\tightlist
\item
  Standardization (Z-score)

  \begin{itemize}
  \tightlist
  \item
    Remove mean
  \item
    Scale to unit variance
  \item
    Properties and use cases
  \end{itemize}
\item
  Min-Max Scaling

  \begin{itemize}
  \tightlist
  \item
    Scale to {[}0,1{]} range
  \item
    When to use
  \item
    Limitations
  \end{itemize}
\item
  Robust Scaling

  \begin{itemize}
  \tightlist
  \item
    Using median and quartiles
  \item
    Handling outliers
  \item
    Trade-offs
  \end{itemize}
\end{enumerate}

\subsection{2. Mathematical Properties of
Scaling}\label{mathematical-properties-of-scaling}

\subsubsection{2.1 Linear Transformations}\label{linear-transformations}

\begin{itemize}
\tightlist
\item
  Scaling as diagonal matrix
\item
  Effect on distances
\item
  Preservation of relationships
\end{itemize}

\subsubsection{2.2 Impact on Statistics}\label{impact-on-statistics}

\begin{itemize}
\tightlist
\item
  Mean and variance
\item
  Correlation coefficients
\item
  Covariance structure
\end{itemize}

\subsubsection{2.3 Numerical
Considerations}\label{numerical-considerations}

\begin{itemize}
\tightlist
\item
  Condition number
\item
  Numerical stability
\item
  Computational efficiency
\end{itemize}

\subsection{3. Measuring Variable
Importance}\label{measuring-variable-importance}

\subsubsection{3.1 Variance-Based Methods}\label{variance-based-methods}

\begin{itemize}
\tightlist
\item
  Total variance contribution
\item
  Explained variance ratio
\item
  Limitations of variance alone
\end{itemize}

\subsubsection{3.2 Correlation Analysis}\label{correlation-analysis}

\begin{itemize}
\tightlist
\item
  Pairwise correlations
\item
  Multiple correlation
\item
  Partial correlation
\end{itemize}

\subsubsection{3.3 Feature Selection}\label{feature-selection}

\begin{itemize}
\tightlist
\item
  Filter methods
\item
  Wrapper methods
\item
  Embedded methods
\end{itemize}

\subsection{4. Practical Applications}\label{practical-applications}

\subsubsection{4.1 When to Scale}\label{when-to-scale}

\begin{itemize}
\tightlist
\item
  Distance-based algorithms
\item
  Gradient-based optimization
\item
  Feature comparison
\end{itemize}

\subsubsection{4.2 When Not to Scale}\label{when-not-to-scale}

\begin{itemize}
\tightlist
\item
  Interpretability requirements
\item
  Natural relationships
\item
  Categorical variables
\end{itemize}

\subsubsection{4.3 Implementation
Considerations}\label{implementation-considerations}

\begin{itemize}
\tightlist
\item
  Pipeline design
\item
  Cross-validation
\item
  Test set handling
\end{itemize}

\subsection{Practice Problems}\label{practice-problems}

\begin{enumerate}
\def\labelenumi{\arabic{enumi}.}
\tightlist
\item
  Scaling Analysis

  \begin{itemize}
  \tightlist
  \item
    Compare different scaling methods
  \item
    Analyze impact on specific algorithms
  \item
    Interpret results
  \end{itemize}
\item
  Variable Importance

  \begin{itemize}
  \tightlist
  \item
    Rank features by importance
  \item
    Compare different metrics
  \item
    Justify choices
  \end{itemize}
\item
  Implementation

  \begin{itemize}
  \tightlist
  \item
    Write scaling functions
  \item
    Handle edge cases
  \item
    Validate results
  \end{itemize}
\end{enumerate}

\subsection{Key Takeaways}\label{key-takeaways}

\begin{enumerate}
\def\labelenumi{\arabic{enumi}.}
\tightlist
\item
  Scaling is crucial for fair comparison
\item
  Different methods suit different needs
\item
  Consider context when choosing method
\item
  Always validate scaling impact
\end{enumerate}

\end{document}
